% Practice-exam template — `exam` document class.
% Toggle the answer key by commenting/uncommenting \printanswers below.
% Build:  bash .claude/skills/practice-exam/build.sh this-file.tex
\documentclass[11pt,addpoints]{exam}

% --- packages (all in TeX Live, verified available) ---
\usepackage[margin=1in]{geometry}
\usepackage{amsmath,amssymb}   % math
\usepackage{booktabs}          % \toprule etc. for tables
\usepackage{array}
\usepackage{enumitem}
\usepackage{graphicx}

% --- answer key: uncomment to render model solutions inline ---
% \printanswers

% --- header / footer ---
\pagestyle{headandfoot}
\firstpageheader{Computational Modelling and Simulation}{}{Practice Exam}
\runningheader{CMAS}{Practice Exam}{\thepage}
\firstpagefooter{}{Page \thepage\ of \numpages}{}
\runningfooter{}{Page \thepage\ of \numpages}{}

\begin{document}

% --- cover page ---
\begin{center}
  {\Large\bfseries The University of Melbourne}\\[2pt]
  School of Computing and Information Systems\\[1em]
  {\Large Computational Modelling and Simulation}\\[4pt]
  {\large Practice Exam}\\[1em]
  Reading time: 15 minutes \quad Writing time: 2 hours
\end{center}
\vspace{1em}
\noindent\textbf{Instructions:} Answer all questions. Marks are shown next to
each part. Total: \numpoints\ marks across \numquestions\ questions.
\vspace{1em}
\hrule
\vspace{1em}

\begin{questions}

% ===== Section A: Short Answer =====
\section*{Section A: Short Answer Questions}

% NOTE: marks go on \part only, never on \question — putting points on both
% makes the `exam` class double-count them in \numpoints. Show the question
% subtotal in the title text instead, as below.
\question
\textbf{General Concepts. [8 marks]}
\begin{parts}
  \part[4] Explain emergence in the context of agent-based modelling, and give
    one example of a globally observable pattern that arises from purely local
    agent rules.
    \begin{solution}
      Emergence is complex, often unpredictable system-level behaviour that
      arises from the interaction of many simple agents following local rules,
      without any agent being explicitly programmed to produce it. Example:
      flocking behaviour (e.g. Boids) emerges from each agent following only
      local alignment, cohesion, and separation rules.
    \end{solution}
  \part[4] What is the difference between a synchronous and an asynchronous
    update scheme in a cellular automaton, and why does the choice matter?
    \begin{solution}
      Synchronous update computes every cell's next state from the current
      grid and applies all changes simultaneously; asynchronous update applies
      changes one cell at a time, so later cells in the same tick may read
      already-updated neighbours. The choice matters because some rule sets
      (e.g. Conway's Game of Life) assume simultaneity and behave differently,
      sometimes unpredictably, under asynchronous update.
    \end{solution}
\end{parts}

% ===== Section B: Method =====
\section*{Section B: Method Questions}

\question
\textbf{Cellular Automata. [12 marks]} Consider a 1D elementary cellular
automaton.
\begin{parts}
  \part[6] Explain how Wolfram's rule numbering scheme (e.g. "Rule 30")
    encodes a 1D CA's update rule, and compute the output bit for the
    neighbourhood \texttt{110} under Rule 30.
    \begin{solution}
      Each of the 8 possible 3-cell neighbourhoods (from \texttt{111} down to
      \texttt{000}) is assigned an output bit; the 8 bits, read as a binary
      number, give the rule number. Rule 30 in binary is \texttt{00011110},
      giving neighbourhood outputs
      $111{\to}0,\,110{\to}0,\,101{\to}0,\,100{\to}1,\,011{\to}1,\,010{\to}1,\,001{\to}1,\,000{\to}0$.
      So \texttt{110} $\to$ \texttt{0}.
    \end{solution}
  \part[6] Why can simple 1D CA rules like Rule 30 or Rule 110 still produce
    highly complex, seemingly random or computationally universal behaviour?
    \begin{solution}
      Even a small local rule space, iterated over many cells and time steps,
      can generate long-range correlations and non-repeating patterns that
      are not predictable from the rule alone without simulation — this is
      Wolfram's notion of computational irreducibility. Rule 110 is proven
      Turing-complete, showing simple local rules can support universal
      computation.
    \end{solution}
\end{parts}

% ===== Section C: Algorithmic =====
\section*{Section C: Algorithmic Questions}

\question
\textbf{Conway's Game of Life. [10 marks]} Consider the 3x3 grid below (1 = live,
0 = dead), using standard Game of Life rules (a live cell survives with 2-3 live
neighbours; a dead cell with exactly 3 live neighbours becomes alive) and an
8-cell (Moore) neighbourhood, ignoring cells outside the grid.
\begin{center}
  \begin{tabular}{ccc}
    \toprule
    0 & 1 & 0 \\
    0 & 1 & 0 \\
    0 & 1 & 0 \\
    \bottomrule
  \end{tabular}
\end{center}
Compute the grid at the next time step and state the resulting pattern's name.
\begin{solution}
  Centre column (blinker): the middle-row cells' neighbour counts change from
  vertical to horizontal alignment. Next state:
  \begin{center}
    \begin{tabular}{ccc}
      \toprule
      0 & 0 & 0 \\
      1 & 1 & 1 \\
      0 & 0 & 0 \\
      \bottomrule
    \end{tabular}
  \end{center}
  This is the classic "blinker" oscillator, alternating between vertical and
  horizontal orientation with period 2.
\end{solution}

\end{questions}
\end{document}
